\section{Introduction}\label{sec:intro}

Deep learning models are increasingly deployed directly on mobile edge devices to deliver low-latency, privacy-preserving, and offline-capable intelligence~\cite{xu2019first,lane2015early}. Consumer smartphone applications routinely embed neural networks for camera image enhancement, real-time object detection, biometric authentication, and natural language understanding~\cite{almeida2021smart,wu2019machine}. However, executing complex deep neural networks within the strict thermal, battery, and memory envelopes of mobile handsets presents acute systems engineering challenges~\cite{lane2016deepx,banbury2021benchmarking}.

Concurrently, multi-platform mobile application development frameworks have achieved widespread industrial adoption~\cite{biornhansen2018survey}. Google's Flutter framework, in particular, has emerged as a dominant stack for modern mobile applications, enabling software teams to target Android and iOS from a single Dart codebase. Unlike legacy hybrid frameworks that rely on interpreted JavaScript bridges or WebView wrappers~\cite{ciman2017empirical,delia2015performance}, Flutter compiles Dart ahead-of-time (AOT) to native ARM machine code and renders interfaces directly via its dedicated graphics engine~\cite{biornhansen2020empirical}.

Despite the widespread adoption of both mobile deep learning and Flutter, a critical architectural gap persists in software engineering practice. On-device inference engines---chiefly Google LiteRT (formerly TensorFlow Lite)---are compiled native C/C++ runtimes. When a Flutter application performs inference, execution must cross a foreign function interface (Dart FFI) boundary. Furthermore, typical production Flutter developers perform data ingestion and preprocessing (such as JPEG image decoding, bilinear resizing, and tensor layout normalization) in pure Dart rather than platform-native code. Developers currently possess no empirical, statistically rigorous evidence quantifying the latency, memory, CPU, and energy costs introduced by this cross-platform software layer. Consequently, architectural decisions between cross-platform and native mobile stacks are guided by anecdotal claims.

Prior mobile machine learning benchmarking efforts, such as AI Benchmark~\cite{ignatov2018ai,ignatov2019ai}, MLPerf Mobile~\cite{reddi2022mlperfmobile}, and library comparative studies~\cite{zhang2022comprehensive,luo2020comparison}, focus almost exclusively on silicon accelerators (CPUs, GPUs, and NPUs) and native C++ runtime engines. These benchmarks evaluate standalone native binaries, deliberately abstracting away the application software frameworks in which deep learning models are actually integrated. Conversely, existing empirical studies of cross-platform performance~\cite{biornhansen2020empirical,willocx2016comparing} evaluate synthetic microbenchmarks or generic UI navigation, omitting high-throughput tensor operations and hardware-delegated inference workloads.

To resolve this empirical blind spot, this paper presents a controlled, full-factorial benchmark study comparing Flutter and Native Android implementations of identical on-device deep learning workloads. By enforcing strict parity---holding model binaries, LiteRT runtime versions, hardware execution delegates, CPU thread configurations, and input sequences identical across both platforms---we isolate and quantify the exact overheads introduced by Flutter.

Our empirical investigation addresses five formal research questions:
\begin{enumerate}
    \item[\textbf{RQ1:}] \textit{Inference Latency.} Is there a statistically or practically significant difference in pure runtime inference latency ($\text{inference\_native\_ms}$) between Flutter and native Android when running identical LiteRT models and delegates?
    \item[\textbf{RQ2:}] \textit{End-to-End Latency Decomposition.} What is the latency discrepancy across the full input-to-result pipeline, and which stage (preprocessing, boundary marshalling, inference, or postprocessing) drives the differential?
    \item[\textbf{RQ3:}] \textit{Resource Utilization.} Does Flutter incur measurable premiums in peak memory consumption (proportional set size, PSS), CPU utilization, or battery drain under continuous inference workloads?
    \item[\textbf{RQ4:}] \textit{Interaction Effects.} How do model scale, precision format (FP32 vs.\ INT8), hardware acceleration delegate (CPU, GPU, NNAPI), and device hardware tiers interact with the cross-platform performance gap?
    \item[\textbf{RQ5:}] \textit{Application Footprint and Startup.} What are the relative penalties in cold-start latency to first inference and packaged application binary size?
\end{enumerate}

This paper provides four primary contributions:
\begin{itemize}
    \item \textbf{C1: Open-Source Parity-Checked Benchmark Suite.} We release open-source Flutter and Android Native benchmark applications implementing an identical execution contract, validated with cryptographic model verification and automated JSONL output validation.
    \item \textbf{C2: Controlled Empirical Characterization.} We execute a multi-tier study across 3 model architectures (MobileNetV3-Large, EfficientDet-Lite0, MobileBERT-QA), 2 numerical precisions (FP32, INT8), 3 execution delegates (CPU/XNNPACK, GPU, NNAPI), and multiple physical smartphone hardware tiers.
    \item \textbf{C3: Pipeline Overhead Decomposition.} We decompose end-to-end latency into data preprocessing, language boundary crossing, kernel invoke execution, and output postprocessing, exposing precisely where execution overheads originate.
    \item \textbf{C4: Evidence-Based Practitioner Guidelines.} We synthesize empirical findings into a structured decision matrix, providing mobile software architects with actionable thresholds for when Flutter is sufficient and when native implementation is justified.
\end{itemize}

The remainder of this paper is structured as follows. Section~\ref{sec:related} reviews related work in mobile deep learning and cross-platform benchmarking. Section~\ref{sec:method} details our experimental methodology, harness contracts, and statistical framework. Section~\ref{sec:results} presents our empirical results corresponding to RQ1--RQ5. Section~\ref{sec:discussion} discusses implications and provides engineering guidelines. Section~\ref{sec:threats} examines threats to validity, and Section~\ref{sec:conclusion} concludes the paper.
